\documentclass[11pt]{article}

\usepackage[review]{acl}
\usepackage{times}
\usepackage{latexsym}
\usepackage{booktabs}
\usepackage{array}
\usepackage[T1]{fontenc}
\usepackage[utf8]{inputenc}
\usepackage{microtype}
\usepackage{inconsolata}
\usepackage{graphicx}

\title{Do Language Models Internalize Reference-Dependent Social Choice?}

\author{First Author \\
  Affiliation / Address line 1 \\
  Affiliation / Address line 2 \\
  \texttt{email@domain} \\\And
  Second Author \\
  Affiliation / Address line 1 \\
  Affiliation / Address line 2 \\
  \texttt{email@domain} \\}

\begin{document}
\maketitle

\begin{abstract}
Large language models (LLMs) are increasingly used as social and economic agents, but it is unclear whether they internalize the structural rules that govern socially dependent choices.
We study this question through the lens of Langtry's model of reference-dependent choice with social comparisons.
The model yields a precise decision rule: an agent combines a private baseline with a closeness-weighted peer reference, while related out-of-domain tasks require matching benefit-cost ratios or resolving relative-status tradeoffs.
We turn these rules into natural-language diagnostic tasks and evaluate six LLMs.
Even the strongest model reaches only 51.6\% on core rule identification; comparative statics are easier but incomplete, social matching remains imperfect, and career sorting is near chance.
Error analysis shows systematic shortcuts, including flat averaging, closest-peer imitation, median anchoring, pure-private responses, and over-socialization in placebo cases.
These diagnostics define the contrastive data needed for targeted fine-tuning before deploying LLMs in simulations of networked social dynamics such as positional asset bubbles.
\end{abstract}

\section{Introduction}

Social comparison is a basic force in economic behavior.
People do not only choose how much to spend, study, work, or invest in isolation; they also respond to what socially salient peers do.
Formal models make this dependence precise by introducing reference points, peer weights, and strategic complementarity into individual utility \cite{clark2008relative, hopkins2004running, langtry2023keeping}.
At the same time, LLMs are increasingly used as simulated social and economic agents \cite{horton2023large, zhou2024real, manning2026general}.
This raises a methodological question for NLP: when an LLM gives a plausible explanation of social pressure, has it learned the structural rule, or merely a surface narrative?

We address this question using Langtry's model of reference-dependent choice with social comparisons \cite{langtry2023keeping}.
The model is attractive as a testbed because it distinguishes several mechanisms that are easy to conflate in text.
The gold decision rule is not simply ``follow the crowd''.
An agent begins from a private baseline and adjusts by a comparison intensity times a closeness-weighted peer reference.
This differs from uniform averaging, nearest-peer imitation, median anchoring, equal mixing, and pure private choice.

We construct a theory-grounded set of natural-language diagnostic tasks.
Each item is generated from latent economic parameters and then rendered as a realistic scenario.
Models are evaluated on rule identification, pairwise comparative statics, placebo controls, social matching, and career sorting.
The same records are designed to support a broader research program: diagnose structural failures, fine-tune models against targeted shortcut contrasts, and then use adapted agents in simulations of social contagion.
This paper reports the diagnostic stage and specifies how the same data support targeted adaptation and simulation.

Our results show that off-the-shelf LLMs are not reliable structural simulators of reference-dependent social choice.
Across six models, Eval-A rule identification ranges from 23.4\% to 51.6\%.
Models perform better on local comparative statics, but still fail on many perturbations involving comparison intensity and peer-weight redistribution.
They also over-socialize placebo decisions and remain near chance on career-status tradeoffs.
The error categories are interpretable, giving concrete targets for subsequent fine-tuning.

\section{Background and Related Work}

\subsection{Decision Making with Social Comparison}
Social comparison has long been used to explain deviations from independent-preference choice.
Classic work on ``keeping up with the Joneses'' and habit formation shows how peer consumption can affect utility and equilibrium outcomes \cite{abel1990asset, gali1994keeping, campbell1999force}.
Status models further connect relative position to risk attitudes, portfolio choices, and visible consumption \cite{robson1992status, bakshi1996spirit, demarzo2004diversification, demarzo2008relative, roussanov2010diversification}.
Empirical and behavioral work documents that relative income, positional goods, and peer effort shape welfare and action choices \cite{clark2008relative, frank1985demand, cohn2014social}.
Langtry's network model is especially useful for our setting because it specifies how peer actions are weighted and yields a compact best-response rule \cite{langtry2023keeping}.

\subsection{LLM-based Social and Market Simulation}
LLMs have recently been studied as simulated economic agents, survey respondents, and participants in social interaction \cite{horton2023large, zhou2024real, manning2026general}.
Related work builds agents for macroeconomic activity, trust behavior, financial decision making, and investor-style tasks \cite{jia2024can, li2024econagent, li2025investorbench}.
These studies show that LLMs can produce behavior that is often plausible at the aggregate level.
Our work asks a complementary question: whether the same models recover the latent rule that generated a decision setting.
This matters because simulation fidelity depends not only on plausible prose, but also on whether micro-level decisions obey the intended behavioral mechanism.

\section{Theory-Guided Task Construction}
\label{sec:tasks}

\subsection{Reference-Dependent Choice}
For individual $i$, let $x_i$ be the action, $\alpha_i$ the comparison intensity, $x_j$ the action of peer $j$, and $g_{ij}$ the normalized closeness weight from $i$ to $j$.
Following Langtry, a baseline utility can be written as
\begin{equation}
u_i =
f\!\left(x_i-\alpha_i\sum_j g_{ij}x_j\right)
-cx_i+b_i\alpha_i\sum_j g_{ij},
\label{eq:utility}
\end{equation}
where $c$ is the marginal cost and $b_i$ captures the value of social connections.
Let
\begin{equation}
R_i=\sum_j g_{ij}x_j
\end{equation}
be the social reference and $F=f'^{-1}(c)$ be the private no-comparison baseline.
The main best response is
\begin{equation}
x_i^*=F+\alpha_iR_i.
\label{eq:main_rule}
\end{equation}
Placebo cases clamp $\alpha_i$ near zero so that $x_i^*\approx F$.

We also include two transfer tasks.
For social-network formation, Langtry's model implies that a person's preferred relationship intensity is governed by their benefit-cost ratio.
We write the social-choice utility as
\begin{equation}
u_i =
f\!\left(x_i-\alpha_i\sum_j g_{ij}x_j\right)
-c_ix_i+b_i\alpha_i\sum_j g_{ij},
\label{eq:social_utility}
\end{equation}
where $b_i$ is the value from social ties and $c_i$ is the marginal cost of the action.
The resulting preferred intensity is
\begin{equation}
x_i^*=\frac{b_i}{c_i}.
\label{eq:social_ratio}
\end{equation}
Pairwise stability therefore requires matched benefit-cost ratios, $b_i/c_i \approx b_j/c_j$.
The social matching task asks which candidate has the closest ratio to the protagonist, using only semantic tiers of benefit and financial constraint in the public prompt.
The career sorting task follows Langtry's firm-choice condition:
\begin{equation}
x_{S,H}-x_{S,L}\geq \alpha_{2i}(\bar{x}_H-\bar{x}_L),
\label{eq:career_threshold}
\end{equation}
where $x_{S,m}$ is the worker's own salary at firm $m$ and $\bar{x}_m$ is the firm average.

\subsection{Natural-Language Tasks}
We sample latent records from the theory and realize them as natural-language scenarios with a teacher model.
The main positional domains cover visible consumption and human-capital effort.
Latents vary along a $3\times3\times3$ grid: comparison intensity $\alpha_i$, closeness-weight dispersion, and peer-action skew.
The teacher uses the structured values to write a realistic scenario, while student-facing prompts hide research notation and require models to infer the rule from the narrative.

We evaluate five splits.
Eval-A asks models to identify which rule best explains a scenario.
Eval-B asks which of two matched scenarios produces a higher action under a single latent perturbation.
The placebo split presents non-positional decisions where peer information is visible but should not affect the answer.
Social OOD asks models to match semantic benefit/cost tiers under Eq.~\ref{eq:social_ratio}.
Career OOD asks models to apply Eq.~\ref{eq:career_threshold} when absolute salary and relative rank conflict.

\section{Experimental Setup}
\label{sec:setup}

\paragraph{Models.}
We evaluate DeepSeek-R1-671B, Qwen, GPT-4, Llama-3.1, Llama-3 8B, and Qwen2 7B.
All models receive the same task prompts and answer-format instructions.

\paragraph{Metric.}
All tasks are multiple-choice and scored by exact parsed answer accuracy.
Unparsed answers are counted as incorrect; parse rates are reported in Appendix~\ref{app:additional_results}.

\paragraph{Evaluation scope.}
This evaluation isolates the diagnostic stage of the broader workflow.
The generated records also support targeted fine-tuning: chosen responses commit to the gold rule, while rejected responses commit to common shortcuts.
After adaptation, we will deploy agents in network simulations of positional asset contagion, where socially visible purchases by central nodes may trigger cascades.

\section{Results}
\label{sec:results}

\begin{table*}[t]
\centering
\small
\setlength{\tabcolsep}{5pt}
\begin{tabular}{@{}lrrrrrr@{}}
\toprule
\textbf{Model} & \textbf{Eval-A} & \textbf{Eval-B} & \textbf{Placebo} & \textbf{Social OOD} & \textbf{Career OOD} & \textbf{Mean} \\
\midrule
DeepSeek-R1-671B & 23.4 & 59.6 & 30.0 & 56.9 & 45.8 & 43.2 \\
Qwen & 33.9 & \textbf{70.0} & 59.7 & 66.0 & 49.5 & 55.8 \\
GPT-4 & \textbf{51.6} & 62.5 & \textbf{96.0} & \textbf{69.4} & 50.0 & \textbf{65.9} \\
Llama-3.1 & 41.7 & 61.8 & 75.2 & 62.0 & 50.3 & 58.2 \\
Llama-3 8B & 40.3 & 58.2 & 79.4 & 59.7 & \textbf{51.7} & 57.9 \\
Qwen2 7B & 46.6 & 56.4 & 91.7 & 59.2 & 47.6 & 60.3 \\
\bottomrule
\end{tabular}
\caption{Diagnostic accuracy (\%). }
\label{tab:main_results}
\end{table*}

\paragraph{Rule identification remains weak.}
Eval-A is the central test of Eq.~\ref{eq:main_rule}.
Even GPT-4 reaches only 51.6\%, and the other models are below 47\%.
This shows that models can read social scenarios without reliably identifying the closeness-weighted rule.

\paragraph{Comparative statics are easier but incomplete.}
Eval-B is consistently higher than Eval-A for most models, with Qwen reaching 70.0\%.
Models are better at judging local directional changes than selecting the full rule.
This partial success is useful for simulation, but insufficient: a simulator must produce the correct action under many network configurations, not only rank two vignettes.

\paragraph{Placebo and transfer expose different failures.}
Placebo accuracy ranges from 30.0\% to 96.0\%, revealing large variation in whether models suppress social comparison when the domain is non-positional.
Social OOD accuracy ranges from 56.9\% to 69.4\%, indicating that semantic benefit-cost matching is nontrivial for current models.
Career OOD remains near chance for all models, suggesting that relative-status tradeoffs are especially brittle when absolute salary and rank conflict.

\section{Error Analysis}
\label{sec:error_analysis}

\begin{table}[t]
\centering
\small
\setlength{\tabcolsep}{3pt}
\begin{tabular}{@{}lrrrr@{}}
\toprule
\textbf{Model} & \textbf{Uniform} & \textbf{Private} & \textbf{Closest/top} & \textbf{Mix} \\
\midrule
DeepSeek-R1 & 21.2 & 8.2 & 26.7 & 17.1 \\
Qwen & 20.5 & 10.3 & 32.2 & 15.5 \\
GPT-4 & 3.6 & 18.4 & 48.5 & 22.3 \\
Llama-3.1 & 7.6 & 27.1 & 33.4 & 12.8 \\
Llama-3 8B & 6.4 & 27.9 & 32.5 & 14.0 \\
Qwen2 7B & 9.5 & 25.7 & 34.0 & 14.0 \\
\bottomrule
\end{tabular}
\caption{Selected Eval-A error categories as a percentage of each model's wrong answers.}
\label{tab:evala_errors}
\end{table}

Errors are structured enough to guide adaptation.
DeepSeek-R1-671B and Qwen often use the wrong aggregator, especially uniform averaging.
GPT-4 and the local models more often collapse the rule to closest-peer or top-peer anchoring; the local models also frequently choose the private baseline in positional settings.
These patterns suggest different fine-tuning contrasts for different models.
For example, uniform-average mistakes require examples where equal averages conflict with closeness weights, while placebo mistakes require examples where peer information is visible but normatively irrelevant.

\section{Toward Adaptation and Simulation}

The diagnostic results determine the next phase of the project.
We will fine-tune selected models with contrast pairs that target the observed shortcut families, then test whether the adapted policy improves both in-distribution rule recovery and transfer to social and career OOD settings.

The final evaluation is a simulation study of positional asset bubbles.
Agents will be placed on a social network and asked whether to buy a socially visible asset, such as a prestige property, limited-edition digital collectible, or meme asset.
We will compare untuned and tuned agents under the same fundamentals and seeded adoption shocks.
Key metrics include adoption rate, cascade depth, sensitivity to central-node seeding, peak willingness-to-pay proxy, and divergence from a no-social-reference baseline.
This setup asks whether micro-level structural adaptation changes macro-level social contagion.

\section{Conclusion}

We introduced a theory-grounded diagnostic suite for reference-dependent social choice and evaluated six LLMs.
The results show a consistent gap between plausible social reasoning and structural rule recovery.
Because the errors are systematic, they provide a principled basis for targeted fine-tuning and for subsequent simulations of positional asset contagion.

\section*{Limitations}

The present version reports only the diagnostic evaluation stage.
Fine-tuning and simulation results remain planned work, so claims about adapted agents are currently methodological rather than empirical.
The data are synthetic and depend on teacher-generated narratives; future versions should include human validation of scenario naturalness and ambiguity.

\bibliography{custom}

\appendix
% Appendix-only float numbering: Table S1, Figure S1, ...
\setcounter{table}{0}
\setcounter{figure}{0}
\renewcommand{\thetable}{S\arabic{table}}
\renewcommand{\thefigure}{S\arabic{figure}}
\section{Additional Diagnostic Results}
\label{app:additional_results}

This appendix provides supporting diagnostics for the main results in
Table~\ref{tab:main_results}. We keep the discussion brief: the purpose is to
make the reported evaluation auditable without shifting the paper's focus away
from the evaluate--adapt--simulate workflow.

\subsection{Evaluation Sizes}

Table~\ref{tab:dataset_sizes} reports the number of evaluated units per split.
Eval-B is counted in paired comparisons rather than individual scenarios,
because each prediction compares two matched records.

\begin{table}[h]
\centering
\small
\begin{tabular}{@{}lr@{}}
\toprule
\textbf{Split} & \textbf{Evaluation units} \\
\midrule
Eval-A rule identification & 1,512 \\
Eval-B comparative statics & 280 pairs \\
Placebo & 504 \\
Social OOD & 576 \\
Career OOD & 576 \\
\bottomrule
\end{tabular}
\caption{Evaluation sizes for the diagnostic stage.}
\label{tab:dataset_sizes}
\end{table}

\subsection{Parse Rates}

Table~\ref{tab:parse_rates} checks whether low accuracy could be explained by
formatting failures. With the exception of GPT-4 on Eval-B, parse rates are
near ceiling. The main conclusions are therefore about decision errors rather
than output-format noncompliance.

\begin{table*}[h]
\centering
\small
\setlength{\tabcolsep}{5pt}
\begin{tabular}{@{}lrrrrr@{}}
\toprule
\textbf{Model} & \textbf{Eval-A} & \textbf{Eval-B} & \textbf{Placebo} & \textbf{Social OOD} & \textbf{Career OOD} \\
\midrule
DeepSeek-R1-671B & 99.9 & 99.3 & 98.4 & 99.8 & 97.0 \\
Qwen & 99.5 & 100.0 & 99.4 & 99.8 & 99.7 \\
GPT-4 & 100.0 & 88.6 & 100.0 & 100.0 & 100.0 \\
Llama-3.1 & 100.0 & 100.0 & 100.0 & 100.0 & 100.0 \\
Llama-3 8B & 100.0 & 100.0 & 100.0 & 100.0 & 100.0 \\
Qwen2 7B & 100.0 & 100.0 & 100.0 & 100.0 & 100.0 \\
\bottomrule
\end{tabular}
\caption{Answer parse rates (\%). Unparsed responses are counted as incorrect in accuracy.}
\label{tab:parse_rates}
\end{table*}

\subsection{Eval-B by Perturbation Type}

Table~\ref{tab:evalb_perturbations} decomposes the comparative-static task by
which latent factor changes. Models are strongest when the peer action level or
weighted reference aggregate increases, but weaker when the perturbation depends
on comparison intensity or the redistribution of closeness weights. This
supports the interpretation that models have partial directional sensitivity but
do not fully internalize the structural rule.

\begin{table*}[h]
\centering
\small
\setlength{\tabcolsep}{4pt}
\begin{tabular}{@{}lrrrrr@{}}
\toprule
\textbf{Model} & \textbf{$F$ up} & \textbf{$\alpha_i$ up} & \textbf{Peer action up} & \textbf{$R_i$ up} & \textbf{Top weight up} \\
\midrule
DeepSeek-R1-671B & 51.8 & 51.8 & 73.2 & 75.0 & 46.4 \\
Qwen & 53.6 & 57.1 & 91.1 & 92.9 & 55.4 \\
GPT-4 & 51.8 & 44.6 & 75.0 & 89.3 & 51.8 \\
Llama-3.1 & 64.3 & 42.9 & 73.2 & 76.8 & 51.8 \\
Llama-3 8B & 53.6 & 55.4 & 67.9 & 66.1 & 48.2 \\
Qwen2 7B & 44.6 & 64.3 & 66.1 & 60.7 & 46.4 \\
\bottomrule
\end{tabular}
\caption{Eval-B accuracy (\%) by comparative-static perturbation.}
\label{tab:evalb_perturbations}
\end{table*}

\subsection{Social OOD by Match Distance}

Table~\ref{tab:social_distance} reports the social OOD split. The task does not
expose the numeric $b/c$ ratio and instead uses semantic benefit/cost tiers. The
mid bucket is consistently hardest, suggesting that models rely on coarse
ordering of semantic cues and struggle when the benefit-cost match is less
separable.

\begin{table*}[h]
\centering
\small
\setlength{\tabcolsep}{6pt}
\begin{tabular}{@{}lrrr@{}}
\toprule
\textbf{Model} & \textbf{Close} & \textbf{Mid} & \textbf{Far} \\
\midrule
DeepSeek-R1-671B & 58.3 & 50.9 & 63.9 \\
Qwen & 67.6 & 60.2 & 72.2 \\
GPT-4 & 74.1 & 60.6 & 75.7 \\
Llama-3.1 & 64.8 & 52.3 & 72.2 \\
Llama-3 8B & 59.7 & 52.8 & 70.1 \\
Qwen2 7B & 60.2 & 51.9 & 68.8 \\
\bottomrule
\end{tabular}
\caption{Social OOD accuracy (\%) by semantic benefit/cost match-distance bucket.}
\label{tab:social_distance}
\end{table*}

\subsection{Eval-A Rule Confusions}

Table~\ref{tab:evala_rule_confusions} expands the Eval-A error analysis. The
dominant mistakes are not random: models repeatedly replace the closeness-
weighted rule with single-peer anchoring, uniform aggregation, pure private
choice, or equal mixing. These categories define the main negative examples for
the targeted fine-tuning stage.

\begin{table*}[h]
\centering
\small
\setlength{\tabcolsep}{3pt}
\begin{tabular}{@{}lrrrrrr@{}}
\toprule
\textbf{Model} & \textbf{Top/closest} & \textbf{Uniform} & \textbf{Private} & \textbf{Median} & \textbf{Counter} & \textbf{Mix} \\
\midrule
DeepSeek-R1-671B & 26.7 & 21.2 & 8.2 & 18.4 & 8.3 & 17.1 \\
Qwen & 32.2 & 20.5 & 10.3 & 17.4 & 3.4 & 15.5 \\
GPT-4 & 48.5 & 3.6 & 18.4 & 6.7 & 0.5 & 22.3 \\
Llama-3.1 & 33.4 & 7.6 & 27.1 & 11.2 & 7.8 & 12.8 \\
Llama-3 8B & 32.5 & 6.4 & 27.9 & 13.1 & 6.1 & 14.0 \\
Qwen2 7B & 34.0 & 9.5 & 25.7 & 15.2 & 1.6 & 14.0 \\
\bottomrule
\end{tabular}
\caption{Eval-A wrong-answer distribution (\% of each model's Eval-A errors).}
\label{tab:evala_rule_confusions}
\end{table*}

\subsection{Career OOD by Status Sensitivity Bucket}

Table~\ref{tab:career_alpha2} stratifies career sorting by relative-status
sensitivity. Accuracy remains near chance across low, mid, and high buckets,
including for GPT-4. This suggests that the failure is not isolated to one
status-sensitivity regime; models generally struggle with the threshold tradeoff
between absolute salary and relative rank.

\begin{table*}[h]
\centering
\small
\setlength{\tabcolsep}{6pt}
\begin{tabular}{@{}lrrr@{}}
\toprule
\textbf{Model} & \textbf{Low $\alpha_{2i}$} & \textbf{Mid $\alpha_{2i}$} & \textbf{High $\alpha_{2i}$} \\
\midrule
DeepSeek-R1-671B & 43.8 & 45.8 & 47.9 \\
Qwen & 47.4 & 49.5 & 51.6 \\
GPT-4 & 50.0 & 50.0 & 50.0 \\
Llama-3.1 & 47.9 & 52.1 & 51.0 \\
Llama-3 8B & 50.5 & 51.6 & 53.1 \\
Qwen2 7B & 47.4 & 46.4 & 49.0 \\
\bottomrule
\end{tabular}
\caption{Career OOD accuracy (\%) by relative-status sensitivity bucket. Accuracy remains close to chance across buckets.}
\label{tab:career_alpha2}
\end{table*}


\end{document}
